\section{Use of AI}
\label{sec:intro-llm-use}

The rapid evolution described in \Cref{sec:intro-motivation} also affected the
practice of my research. I now use two LLM-based assistants as part of my
research workflow: ChatGPT~\cite{openai2022chatgpt} and
Claude~\cite{anthropic2023claude}. I introduced them into this workflow only in
the third year of the PhD, when their capabilities became sufficient for
sustained technical assistance. I use them for brainstorming, program
exploration, debugging and refactoring, Lean proof search, and the revision of
technical prose. I also use them as critics, asking them to identify unclear
explanations, propose alternative proof strategies, and point to cases that
merit further investigation.

I remain responsible for validating every output before incorporating it into
the thesis or its software artefacts. This includes checking formal developments
with Lean, exercising software changes through the relevant builds and tests,
verifying factual and bibliographic claims against primary sources, and
reviewing suggested statements and models against the intended B semantics. The
final selection, interpretation, and presentation of all material are mine.
